% !PRINTED 154
\weberpage{154}

Erhebt man ins Quadrat, so erhält man für die Discriminante der cubischen Gleichung (1) nach \S.\,35, (6)
\weq{6}{D = -(27 b^2 + 4 a^3).}
Hieraus erhält man die Discriminante der allgemeinen Function
\weq{7}{a_0 x^3 + a_1 x^2 + a_2 x + a_3,}
wenn man für $a$ und $b$ die Ausdrücke \S.\,34, (4) einsetzt, in denen man $a_1$, $a_2$, $a_3$ durch
\[
\frac{a_1}{a_0}, \quad \frac{a_2}{a_0}, \quad \frac{a_3}{a_0}
\]
ersetzt und [\S.\,46, (3)] mit $a_0^4$ multiplicirt. Man erhält so
\[
\begin{aligned}
a &= \frac{3 a_0 a_2 - a_1^2}{3 a_0^2},\\
b &= \frac{2 a_1^3 - 9 a_0 a_1 a_2 + 27 a_0^2 a_3}{27 a_0^3},
\end{aligned}
\]
und folglich
\weq{8}{D = -\frac{1}{27 a_0^2}\left\{(2 a_1^3 - 9 a_0 a_1 a_2 + 27 a_0^2 a_3)^2 + 4(3 a_0 a_2 - a_1^2)^3\right\}.}

Wenn man das Quadrat und den Cubus ausrechnet, so kommt man zu der Formel (10) des vorigen Paragraphen zurück; es ist aber bemerkenswerth, dass die Formel (8) scheinbar den Nenner $27 a_0^2$ enthält, der sich beim Ausrechnen forthebt.

Ebenso können wir bei den biquadratischen Gleichungen verfahren. Wir erhalten, wenn $\alpha$, $\beta$, $\gamma$, $\delta$ die Wurzeln der biquadratischen Gleichung
\weq{9}{x^4 + a x^2 + b x + c = 0}
sind, nach \S.\,37, (5)
\weq{10}{\begin{aligned}
\alpha - \beta &= v + w, & \gamma - \delta &= v - w,\\
\alpha - \gamma &= w + u, & \delta - \beta &= w - u,\\
\alpha - \delta &= u + v, & \beta - \gamma &= u - v.
\end{aligned}}
Danach wird die Discriminante der Gleichung (9)
\[
D = (v^2 - w^2)^2 (w^2 - u^2)^2 (u^2 - v^2)^2,
\]
d. h. es ist $D$ gleich der Discriminante der cubischen Resolvente \S.\,37, (4)
\weq{11}{z^3 + 2 a z^2 + (a^2 - 4c) z - b^2 = 0,}
deren Wurzeln $u^2$, $v^2$, $w^2$ sind.

Wenn man diese Discriminante aber nach der Formel (10), \S.\,46 bildet, indem man
%CONT>>
